\documentclass[12pt]{article}
\usepackage[margin=1in]{geometry}
\usepackage{amsmath, amssymb, mathtools}
\usepackage{enumitem}
\usepackage{booktabs}
\usepackage{hyperref}
\hypersetup{colorlinks=true,linkcolor=black,urlcolor=blue,citecolor=blue}

\newcommand{\R}{\mathbb{R}}
\newcommand{\dd}{\,\mathrm{d}}
\newcommand{\dtotal}[2]{\frac{\mathrm{d} #1}{\mathrm{d} #2}}
\newcommand{\pd}[2]{\frac{\partial #1}{\partial #2}}
\setlist[enumerate]{itemsep=0.8em}
%\setlist[enumerate,1]{label=\textbf{Q\arabic*.}}

\title{Practice Problem Set 1}
\author{ECON 320 -- Introduction to Mathematical Economics}
\date{Fall 2025}

\begin{document}
\maketitle

\noindent\textbf{Instructions.} Show all essential steps. Unless specified, provide exact answers (no decimals). Clearly state units and interpret results when an economic context is given.


\begin{enumerate}
\item Solve each inequality for $x$ and express the solution in interval notation.
\begin{enumerate}[label=(\alph*)]
  \item $\lvert 2x-3\rvert \le 5$.
  \item $(x-3)(x+2)\le 0$.
  \item $\displaystyle \frac{1}{x-2}\ge 3$.
\end{enumerate}

\item Let $A=[-1,4)$ and $B=(1,6]$. Find:
\begin{enumerate}[label=(\alph*)]
  \item $A\cup B$ and $A\cap B$,
  \item $A\setminus B$ and $B\setminus A$,
  \item $\inf(A\cap B)$, $\sup(A\cap B)$.
\end{enumerate}

\item Determine the \emph{domain} of
\[
f(x)=\frac{\sqrt{\,5-\ln(x^{2}-9)\,}}{x-3}.
\]
Write the domain as a union of intervals and justify briefly.

\item Consider $f(x)=\dfrac{2x-3}{x+1}$ with domain $x>-1$, $x\neq -1$.
\begin{enumerate}[label=(\alph*)]
  \item Find the inverse function $f^{-1}(y)$ and state its domain.
  \item Evaluate $f^{-1}(5)$.
\end{enumerate}

\item Exponential--logarithmic equations. Solve for $x$ (exact form preferred):
\begin{enumerate}[label=(\alph*)]
  \item $e^{3x}=7$.
  \item $2^{x+1}=5e^{0.2x}$ (give $x$ to three decimal places).
\end{enumerate}

\item Sequences and bounds. Let $S=\left\{\dfrac{2n+1}{n+3}: n\in\mathbb{N}\right\}$.
\begin{enumerate}[label=(\alph*)]
  \item Compute $\displaystyle \lim_{n\to\infty} \frac{2n+1}{n+3}$.
  \item Find $\inf S$ and $\sup S$ (justify).
\end{enumerate}

\item Sets in $\R^2$. Let $B=\{(x,y)\in\R_{\ge 0}^2: 2x+3y\le 60\}$.
\begin{enumerate}[label=(\alph*)]
  \item Find the $x$- and $y$-intercepts of the boundary $2x+3y=60$.
  \item Compute the area of $B$.
  \item Write $B$ as $\{(x,y): 0\le x\le \bar x,\; 0\le y\le g(x)\}$ by solving for $y$ in terms of $x$.
\end{enumerate}

\item Composition and transformation. Let $u(x)=\ln x$ (domain $x>0$) and $g(x)=2-3\,\ln(4-x)$.
\begin{enumerate}[label=(\alph*)]
  \item State the domain of $g$.
  \item Write $g$ as an affine transformation of $u\!\left(\cdot\right)$ composed with a linear function. Identify the horizontal shift and reflection.
\end{enumerate}

\item Absolute value with parameters. For $a\in\R$ and $b>0$, solve
\[
\lvert x-a\rvert \le b,
\]
and express the solution set as an interval in terms of $a$ and $b$.

\item Domains with radicals and rationals. Determine the domain of
\[
h(x)=\sqrt{9-x^2}\;+\;\frac{1}{\sqrt{x-1}}.
\]
Write the final domain as a (possibly empty) interval or union of intervals.
\end{enumerate}

\begin{enumerate}[resume]
\item Differentiate and simplify. Then evaluate at the indicated point.
\begin{enumerate}[label=(\alph*)]
  \item $y=\dfrac{(3x^2+1)^5}{x\,e^{2x}}$, find $y'(1)$.
  \item $y=\ln\!\big(\sqrt{1+4x^2}\big)$, find $y'$ and $y''$.
\end{enumerate}

\item Cost and average cost. Let $C(q)=0.5q^2+10q+400$ with $q\ge 0$.
\begin{enumerate}[label=(\alph*)]
  \item Find marginal cost $MC(q)$ and average cost $AC(q)$.
  \item Compute $MC(20)$ and $AC(20)$.
  \item Find the output $\hat q$ that minimizes $AC(q)$ and the minimum average cost.
\end{enumerate}

\item Elasticity of demand.
\begin{enumerate}[label=(\alph*)]
  \item For $Q(P)=1000P^{-1.2}$, compute the price elasticity $\varepsilon(P)$ and state whether it depends on $P$.
  \item For $Q(P)=200-2P$, compute $\varepsilon$ at $P=50$.
\end{enumerate}

\item Implicit differentiation. For the curve
\[
x e^{x}+y^2=10,
\]
find $\dfrac{\mathrm{d}y}{\mathrm{d}x}$ as a function of $(x,y)$ and then evaluate at the point with $x=1$ and $y>0$.

\item Total differential and proportional changes. Suppose output is
\[
Q=K^{0.3}L^{0.7}P^{0.1}.
\]
At the base point $(K,L,P)=(100,400,2)$, approximate the percentage change in $Q$ if $K$ increases by $5\%$, $L$ decreases by $2\%$, and $P$ increases by $10\%$ (use differentials).

\item Linear supply--demand with a per-unit tax $t$. Let
\[
Q_s=-20+4(P-t),\qquad Q_d=500-6P+2Y,
\]
with $Y=100$.
\begin{enumerate}[label=(\alph*)]
  \item Solve for the equilibrium price $P(t)$ and quantity $Q(t)$ as functions of $t$.
  \item Compute $\dfrac{\mathrm{d}P}{\mathrm{d}t}$ and $\dfrac{\mathrm{d}Q}{\mathrm{d}t}$ and evaluate at $t=5$.
  \item Briefly interpret $\dfrac{\mathrm{d}P}{\mathrm{d}t}$ in words (tax incidence on buyers).
\end{enumerate}

\item Profit maximization and comparative statics. Let $\pi(q)=pq-(aq+\tfrac{b}{2}q^2)$ with parameters $p>0$, $b>0$.
\begin{enumerate}[label=(\alph*)]
  \item Find the profit-maximizing $q^*(p,a,b)$.
  \item Compute $\pd{q^*}{p}$ and $\pd{q^*}{a}$ and sign them.
\end{enumerate}

\item Elasticity of supply. If $Q(P)=cP^{\eta}$ with $c>0$, $\eta>0$,
\begin{enumerate}[label=(\alph*)]
  \item Derive the price elasticity of supply $\varepsilon_s(P)$.
  \item For $c=5$ and $\eta=0.8$, compute $\varepsilon_s$ at $P=4$ and $P=25$ and comment on constancy.
\end{enumerate}

\item Implicit-function comparative statics. Suppose $F(x,a)=x^3-4x+a=0$ defines $x=x(a)$ near a root.
\begin{enumerate}[label=(\alph*)]
  \item Use the implicit-function rule to find $\dfrac{\mathrm{d}x}{\mathrm{d}a}$ in terms of $x$.
  \item Evaluate $\dfrac{\mathrm{d}x}{\mathrm{d}a}$ at a root $x=2$.
\end{enumerate}

\item Multivariate total derivative (dynamic comparative statics). Let market equilibrium be defined by
\[
\begin{cases}
D(P,Y)=\alpha Y - \beta P = Q,\\[0.3em]
S(P,W)=\gamma P - \delta W = Q,
\end{cases}
\]
with positive parameters. Treat $Y=Y(t)$ and $W=W(t)$ as time-varying exogenous variables.
\begin{enumerate}[label=(\alph*)]
  \item Solve for $P$ and $Q$ as functions of $Y$ and $W$.
  \item Compute $\dtotal{P}{t}$ and $\dtotal{Q}{t}$ in terms of $\dot Y$ and $\dot W$.
  \item Evaluate $\dtotal{Q}{t}$ when $\alpha=\beta=\gamma=\delta=2$, $Y=100$, $W=60$, $\dot Y=3$, $\dot W=-1$.
\end{enumerate}

\item Log-linear approximation. Let $R(q)=Aq^\theta$ with $A>0$, $\theta\in(0,1)$ and cost $C(q)=c_0+c_1 q$ ($c_0,c_1>0$). Profit $\pi(q)=R(q)-C(q)$.
\begin{enumerate}[label=(\alph*)]
  \item Compute $MR(q)$ and $MC(q)$.
  \item Using a first-order (log) approximation around $q=q_0$, derive an approximate expression for the change in profit $\Delta\pi$ from a small percentage change in $q$.
  \item At $A=200$, $\theta=0.6$, $c_0=100$, $c_1=20$, $q_0=50$, estimate $\Delta\pi$ for a $+2\%$ change in $q$.
\end{enumerate}

\item Comparative statics with a subsidy. Demand: $Q_d=300-3P$. Supply with per-unit subsidy $s>0$: $Q_s=-30+2(P+s)$.
\begin{enumerate}[label=(\alph*)]
  \item Solve for the equilibrium $P(s)$ and $Q(s)$.
  \item Compute $\dfrac{\mathrm{d}P}{\mathrm{d}s}$ and $\dfrac{\mathrm{d}Q}{\mathrm{d}s}$ and interpret the signs.
\end{enumerate}



\item Implicit elasticity. Suppose demand satisfies $F(Q,P)=\ln Q+\phi Q-\ln P=0$ with $\phi>0$ (defining $Q=Q(P)$).
\begin{enumerate}[label=(\alph*)]
  \item Using implicit differentiation, find $\dfrac{\mathrm{d}Q}{\mathrm{d}P}$.
  \item Derive the price elasticity $\varepsilon=\dfrac{\mathrm{d}Q}{\mathrm{d}P}\cdot \dfrac{P}{Q}$ in terms of $Q$ (or $P$).
\end{enumerate}



\end{enumerate}

\end{document}
